% AulaTeX - maqueta inicial de construcción descendente
% Contrato futuro: el Agente puede investigar, redactar, evaluar y compilar a partir de este archivo.
\documentclass[12pt]{article}
\usepackage[utf8]{inputenc}
\usepackage[T1]{fontenc}
\usepackage[spanish]{babel}
\usepackage{enumitem}
\usepackage{hyperref}
\usepackage{longtable}

\title{UAS}
\author{AulaTeX}
\date{\today}

\begin{document}
\maketitle

\noindent\textbf{Nodo editorial:} UAS\\
\textbf{Padre editorial:} interinstitucional\\
\textbf{Modo de generación:} reforzar\\
\textbf{Destino:} UAS\\
\textbf{Contrato futuro:} ready-for-agent

\section*{Ingesta base}
\begin{itemize}[leftmargin=*]
  \item Texto libre proporcionado: sí
  \item Documento de apoyo proporcionado: no
\end{itemize}

\subsection*{Resumen de la ingesta textual}
\begin{quote}
https://www.uas.edu.mx/
\end{quote}

\section*{Objetivo}
\begin{itemize}[leftmargin=*]
  \item Establecer una base institucional compilable y reutilizable para producción académica en LaTeX.
  \item Estandarizar entradas de reporte, actividad y presentación para programas UAS.
  \item Completar la maqueta inicial de UAS.
  \item Disponer de una base institucional compilable para producción académica UAS.
  \item Estandarizar estructura mínima de reporte, actividad y presentación.
  \item Facilitar escalamiento por carrera/materia con trazabilidad editorial.
\end{itemize}

\section*{Competencias}
\begin{itemize}[leftmargin=*]
  \item Gestión editorial académica con estructura modular.
  \item Uso de plantillas LaTeX institucionales con control bibliográfico.
  \item Trazabilidad entre investigación, redacción y compilación.
  \item Gestión de documentos académicos LaTeX bajo contrato institucional.
  \item Aplicación de criterios de calidad editorial y citación verificable.
  \item Organización modular de contenidos por carrera y materia.
\end{itemize}

\section*{Resultados esperados}
\begin{itemize}[leftmargin=*]
  \item Nodo UAS con artefactos canónicos listos para completar.
  \item Rutas y convenciones estables para crecimiento por carrera/materia.
  \item Checklist mínimo de calidad para compilación y citación.
  \item Entradas .tex institucionales listas para completar por agente.
  \item Estructura de carpetas estable y reutilizable.
  \item Bibliografía base institucional preparada para crecimiento controlado.
\end{itemize}

\section*{Estructura sugerida}
\begin{itemize}[leftmargin=*]
  \item Portada institucional y metadatos de asignatura.
  \item Objetivo y competencias de la actividad/reporte.
  \item Desarrollo por secciones con evidencias y citas.
  \item Criterios de evaluación/rúbrica resumida.
  \item Conclusiones, referencias y anexos (si aplica).
  \item Portada y metadatos institucionales/materia.
  \item Objetivo, competencias y alcance de la entrega.
  \item Desarrollo seccionado con evidencias y citas.
  \item Criterios de evaluación o rúbrica breve.
  \item Cierre, referencias y anexos opcionales.
  \item Bloque de pendientes de investigación en comentarios.
\end{itemize}

\section*{Criterios de evaluación}
\begin{itemize}[leftmargin=*]
  \item Pertinencia académica respecto a programa/materia validada.
  \item Claridad argumentativa y coherencia interna.
  \item Uso correcto de citas y referencias verificables.
  \item Compilación limpia y formato consistente.
  \item Coherencia entre propósito, desarrollo y evaluación.
  \item Cumplimiento de formato y compilación sin errores.
  \item Uso correcto de fuentes verificables y citas trazables.
  \item Consistencia terminológica en toda la entrega.
\end{itemize}

\section*{Bibliografía requerida}
\begin{itemize}[leftmargin=*]
  \item Archivo base: bibliografia-uas.bib.
  \item Fuentes oficiales UAS y documentos académicos verificables.
  \item Norma de citación definida por plantilla o lineamiento de curso.
  \item bibliografia-uas.bib como base institucional.
  \item Fuentes oficiales UAS y normativa académica pública.
  \item Fuentes disciplinares validadas en fase de investigación.
\end{itemize}

\section*{Marcadores de investigación}
\begin{itemize}[leftmargin=*]
  \item Refuerzo Ciclo A materializado
  \item Refuerzo Ciclo A materializado
  \item Refuerzo Ciclo A materializado
  \item Refuerzo Ciclo A materializado
  \item Definir fuentes, citas y vacíos de contenido antes de redactar.
  \item Refuerzo Ciclo A materializado catálogo oficial de carreras UAS vigente
  \item Refuerzo Ciclo A materializado programas analíticos por materia prioritaria
  \item Refuerzo Ciclo A materializado lineamientos de evaluación institucionales o por unidad académica
  \item Refuerzo Ciclo A materializado corpus bibliográfico troncal por disciplina
\end{itemize}

\section*{Indicaciones editoriales por entregable}

\subsection*{Plantilla}
\begin{itemize}[leftmargin=*]
  \item \texttt{\% Archivo: UAS/reporte-uas.tex}
  \item \texttt{\% Estructura mínima: input\{template\} + metadatos + secciones base + bibliography\{bibliografia-uas\}}
  \item \texttt{\% Mantener compatibilidad con TEXINPUTS/BIBINPUTS del repositorio.}
  \item Definir plantilla base con reglas editoriales, assets y referencias de estilo antes de redactar.
  \item \texttt{\% UAS/reporte-uas.tex}
  \item \texttt{\% input\{template\} desde base/ por TEXINPUTS}
  \item \texttt{\% Metadatos: institucion, carrera, materia, periodo, autor}
  \item \texttt{\% Secciones mínimas: objetivo, competencias, desarrollo, evaluacion, cierre}
  \item \texttt{\% bibliography\{bibliografia-uas\}}
  \item Prepara una raíz institucional reutilizable con carpetas por carrera, carpeta assets y referencias claras a plantillas LaTeX compartidas.
  \item Usar punto de entrada .tex por tipo de entrega: reporte, actividad, presentacion.
  \item Resolver plantillas mediante TEXINPUTS configurado en .latexmkrc; no hardcodear rutas absolutas.
  \item Resolver .bib mediante BIBINPUTS y nombre consistente por institución/materia.
  \item Conservar raíz con archivos canónicos institucionales en UAS/.
  \item Crear/normalizar carpeta assets/ para recursos gráficos y tablas reutilizables.
  \item Organizar contenido académico por carrera/programa dentro de subcarpetas dedicadas.
  \item Usa la ingesta proporcionada como restricción editorial inicial.
\end{itemize}

\subsection*{Actividad}
\begin{itemize}[leftmargin=*]
  \item \texttt{\% Archivo sugerido por materia: UAS/<carrera>/<materia>/actividad-<slug>.tex}
  \item \texttt{\% Secciones: contexto, instrucciones, desarrollo, evidencia, criterios de evaluación, referencias.}
  \item \texttt{\% Incluir marcador de investigación pendiente en comentarios.}
  \item Desarrollar la actividad respetando objetivo, criterios de evaluación y marcadores de investigación.
  \item \texttt{\% UAS/<carrera>/<materia>/actividad-<slug>.tex}
  \item \texttt{\% Secciones: contexto, instrucciones, entregables, criterios, referencias}
  \item \texttt{\% Incluir bloque Refuerzo Ciclo A materializado cuando aplique}
  \item \texttt{\% Mantener actividad como plantilla incompleta, no contenido final}
  \item Establecer una base institucional compilable y reutilizable para producción académica en LaTeX.
  \item Estandarizar entradas de reporte, actividad y presentación para programas UAS.
  \item Pertinencia académica respecto a programa/materia validada.
  \item Claridad argumentativa y coherencia interna.
  \item Uso correcto de citas y referencias verificables.
  \item Validar oferta académica y denominaciones oficiales en https://www.uas.edu.mx/.
  \item Confirmar estructura de carreras/facultades para crear subcarpetas reales.
  \item Recuperar lineamientos de evaluación o formato institucional, si existen públicamente.
  \item Usa la ingesta proporcionada como restricción editorial inicial.
\end{itemize}

\subsection*{Reporte}
\begin{itemize}[leftmargin=*]
  \item \texttt{\% Archivo sugerido por materia: UAS/<carrera>/<materia>/reporte-<slug>.tex}
  \item \texttt{\% Secciones: introducción, objetivos, marco conceptual, desarrollo, conclusiones, bibliografía.}
  \item \texttt{\% Usar bibliografía institucional o específica de materia según volumen.}
  \item Estructurar el reporte con bibliografía requerida y cierre verificable.
  \item \texttt{\% UAS/<carrera>/<materia>/reporte-<slug>.tex}
  \item \texttt{\% Secciones: introduccion, objetivos, marco, desarrollo, conclusiones, referencias}
  \item \texttt{\% Reusar estructura institucional y ajustar a materia}
  \item \texttt{\% Citar solo entradas existentes en .bib verificado}
  \item UAS/reporte-uas.tex
  \item UAS/presentacion-uas.tex
  \item UAS/bibliografia-uas.bib
  \item Redacción formal académica en español claro, sin tono promocional.
  \item Bullets breves y accionables en guías editoriales.
  \item Consistencia terminológica entre portada, objetivos, competencias y evaluación.
  \item Priorizar archivo institucional bibliografia-uas.bib como contenedor maestro inicial.
  \item Permitir .bib por materia cuando el volumen crezca, manteniendo trazabilidad.
  \item Registrar solo fuentes verificables y citables; sin placeholders falsos.
  \item Compila sin errores con scripts/latexmk-build.ps1 pasando solo el .tex objetivo.
  \item Cada .tex declara bibliografía institucional o de materia existente.
  \item No hay rutas rotas de assets ni includes de plantillas.
\end{itemize}

\subsection*{Presentación}
\begin{itemize}[leftmargin=*]
  \item \texttt{\% Archivo: UAS/presentacion-uas.tex o por materia.}
  \item \texttt{\% Estructura: portada, objetivo, puntos clave, cierre, referencias.}
  \item \texttt{\% Diseñar para síntesis visual; contenido extenso queda en reporte.}
  \item Preparar la presentación con síntesis visual, continuidad editorial y evidencias clave.
  \item \texttt{\% UAS/presentacion-uas.tex o UAS/<carrera>/<materia>/presentacion-<slug>.tex}
  \item \texttt{\% Estructura: portada, objetivo, hallazgos clave, cierre, referencias}
  \item \texttt{\% Priorizar sintesis visual y consistencia con reporte asociado}
  \item \texttt{\% Sin duplicar desarrollo extenso del reporte}
  \item Nodo UAS con artefactos canónicos listos para completar.
  \item Rutas y convenciones estables para crecimiento por carrera/materia.
  \item Checklist mínimo de calidad para compilación y citación.
  \item Redacción formal académica en español claro, sin tono promocional.
  \item Bullets breves y accionables en guías editoriales.
  \item Reforzamiento de memoria, estructura y convenciones de archivos.
  \item Definición de maqueta inicial reusable para actividades y reportes.
  \item Sintetiza visualmente la memoria editorial sin perder continuidad con plantilla, actividad y reporte.
\end{itemize}

% --- Ciclo A: refuerzo editorial materializado ---
\section*{Refuerzo editorial Ciclo A}
Este apartado consolida la mejora editorial incremental del nodo a partir del estándar maduro de AulaTeX. El refuerzo atiende brechas detectadas de bibliografía, citas, análisis propio, transferencia y cierre argumentado, sin sustituir la consigna original ni copiar contenido de los casos maestros.

\subsection*{Tesis de trabajo}
La actividad debe sostener una postura verificable: el producto solicitado no sólo organiza información, sino que permite interpretar el problema disciplinar, justificar decisiones y reconocer sus consecuencias académicas o profesionales.

\subsection*{Análisis propio}
El hallazgo central es que una entrega madura requiere pasar de la descripción a la interpretación. Por ello, el desarrollo debe explicar qué revela el producto construido, por qué es relevante para la materia y qué criterio permite evaluar su pertinencia. Esta operación convierte la evidencia reunida en argumento académico.

\subsection*{Transferencia profesional}
La transferencia consiste en vincular el producto con situaciones reales de desempeño: diagnóstico, toma de decisiones, comunicación de resultados, cumplimiento normativo, intervención educativa o resolución de problemas según el campo disciplinar. Así, la actividad adquiere valor formativo más allá de la plantilla.

\subsection*{Cierre argumentado}
En consecuencia, la calidad de la actividad depende de que el producto visible, las fuentes y la conclusión formen una secuencia coherente. La conclusión debe recuperar la tesis, indicar el principal aprendizaje y señalar una implicación práctica o conceptual.

% Brechas locales atendidas por Ciclo A: tesis, analisis_propio, conclusion_argumentada, bibliografia_insuficiente, citas_insuficientes, pendientes_editoriales
% Reglas globales aplicadas:
% - Priorizar cierre de brecha recurrente: bibliografia_insuficiente (427 nodos).
% - Priorizar cierre de brecha recurrente: conclusion_argumentada (375 nodos).
% - Priorizar cierre de brecha recurrente: citas_insuficientes (328 nodos).
% - Priorizar cierre de brecha recurrente: analisis_propio (288 nodos).
% --- Fin Ciclo A ---


\subsection*{Citas de refuerzo Ciclo A}
La mejora editorial se vincula con la memoria documental disponible mediante las referencias \cite{cicloARefuerzo01,cicloARefuerzo02,cicloARefuerzo03,cicloARefuerzo04,cicloARefuerzo05,cicloARefuerzo06}. Estas citas deben revisarse en una etapa disciplinar fina para sustituir o complementar fuentes genéricas por literatura específica del tema.

\end{document}
